% !TeX root = ../main.tex

% DOCX body[169]
\section{Conclusion}\label{sec:conclusion}

% DOCX body[170]
In this paper, we proposed SGTG, a semantic-guided Gaussian temporal grounding framework for audio-visual generalized zero-shot learning. SGTG addresses insufficient discrimination between similar class embeddings, dilution of local discriminative cues by redundant background content, and modality interference under weak audio-visual correspondence. First, fine-grained visual and auditory descriptions, together with CLIP and CLAP text encoders, construct class-semantic representations. DSPO then improves class separability while preserving semantic relationships, providing reliable semantic references for temporal evidence extraction and unseen-class recognition. Second, CGTA incorporates class semantics into audio-visual sequence encoding. Independent visual and audio Gaussian experts with adaptive routing continuously localize, weight, and aggregate class-relevant segments, reducing the dilution of local evidence by irrelevant backgrounds. Finally, DMIF models intra-modal and inter-modal relationships separately and dynamically combines intra-modal representations, inter-modal representations, and direct temporal evidence according to the sample and candidate class. This design preserves complementary information while mitigating modality interference and fusion bias. Experiments demonstrate strong performance on VGGSound-GZSL\textsuperscript{\textit{cls}}, UCF-GZSL\textsuperscript{\textit{cls}}, and ActivityNet-GZSL\textsuperscript{\textit{cls}}. Future work will investigate lightweight class-conditioned temporal modeling and noise-robust fusion to improve applicability to large-scale audio-visual generalized zero-shot recognition and complex audio-visual scenes.
